
\paragraph{}In order to make lateral connections parallelisable, the SSM framework requires the constraint of linear recurrence. Applying the same constraint to feedback connections implies making the feed forward pathway linear as well, which is quite limiting. Indeed, with linear feedforwards and linear recurrence, the model becomes entirely linear and can be replaced by a single linear layer. A natural way to alleviate this issue is through linear skip connections. This creates many parallel feedforward pathways, one of which involves only the linear skip connections. The feedback pathway can then project from this linear skip connection pathway back to the lower layer, making it linear and parallelisable, while maintaining non linearity in the overall feedforward. This whole feedback pathway can then be collapsed into a single linear recurrence exactly as in the SSM framework. This idea is illustrated in \cref{fig:feedback_connections}.
%
Thus, linear recurrent layers give at once efficient parallelisable lateral connections, and efficient parallelisable feedback connections under the constraint that the feedback projects from linear skip connections. Parallelisable non linear recurrence and feedback are thus nothing but the very same problem, and solving this problem is an exciting avenue for future work.

In the context of biological neural circuits, this constraint finds a natural conterpart. As mentioned in \cref{sec:neural_circuits_brain_science}, feedback connections are known to typically follow fast magnocellular pathways, while full feedforward pathways typically include slower parvocellular pathways. Thus, the constraint that feedback connections project from linear skip connections aligns well with the known biology of the visual system, providing a biologically plausible framework for modeling feedback in neural circuits. \greg{the above might be too much ?} Additionally, in the visual cortex, feedback connections provide rapid, long range but coarse top-down modulation, while lateral connections are finer but slower and locally constrained. This finds a perfect equivalent in our convolutional framework, where feedback connections can be modelled as convolutions with large kernels, while lateral connections are modelled as convolutions with small kernels. Large kernels are typically computationally expensive, but since we work in spectral domain, they come with zero additional cost compared to smaller ones, and even allow for more efficient parameterisation directly in spectral domain without the need to go back and forth to spatial domain to impose kernel size constraints. Small kernels for lateral connections can also be modeled, and separately from their larger counterparts, allowing for more efficient parameterisation and interpretability. Indeed, \citet{TODO} describes a way to parameterise orthogonal convolutions in spatial domain, allowing for a parameterisation of our discretized operator (\Cref{eq:conv_ssm_discrete:1}). \greg{Future work is needed to (i) make this parameterisation more expressive (cf discussion with thibaut), (ii) parameterise the continuous time operator, which does not need to be orthogonal but to have eigenvalues with negative real part for stability.}

For biological plausibility or for interpretability sake, these two parameterisations can be separated into two consecutive ConvSSM layers, though this comes at some computational cost as both layers could be merged into a single one for efficiency. \greg{TODO : make ablation experiment to see if we gain in practice from separating lateral and feedback connections or if this does not help at all and a single "feedback" layer is enough.}
